\infographicchapter{Fundus Photography}{fundus.png}
\label{ch:infographic-fundus}

\section{Overview}

Fundus photography is a quick, painless, noninvasive eye exam that takes crisp color pictures of the back inside surface of your eye (the fundus) using a specialized camera. The camera shines gentle light through the outer ring of your pupil and captures the reflected image through the center of your pupil, neatly avoiding annoying surface glare. Standard cameras photograph the central 30--50 degrees of the retina, while newer ultra-widefield systems can capture up to 140--200 degrees---seeing nearly all the way to the far edges of the retina in a single shot.

Fundus photos clearly show the retina (the light-sensing lining of the eye), the optic nerve head, the center of sharp vision (macula and fovea), and the tiny retinal blood vessels. Eye doctors use these photos every day to screen for and track diabetic eye disease (diabetic retinopathy), optic nerve changes in glaucoma, age-related macular degeneration (AMD), and blocked retinal blood vessels. They also help document eye inflammation, watch harmless eye freckles (nevi) or eye tumors for any growth, support remote screening clinics, and keep a permanent baseline record of your eye health.

The biggest benefits of fundus photography are that it is fast, painless, repeatable, and uses zero X-rays and no dye injections. Its main limits are that cloudy lenses (cataracts) or bleeding inside the eye can blur the photos, small pupils can narrow the view unless dilating drops are used, and the pictures are flat two-dimensional views that cannot measure fluid depth or watch live dye leakage. Overall safety is outstanding, with only a brief bright camera flash and, when dilating eye drops are used, temporary blurry near vision and a rare risk of an eye-pressure spike in narrow-angle eyes.

\section{Underlying Physical Principles}

\begin{itemize}
    \item \textbf{Glare-Free Ring Lighting:} The fundus camera works like a gentle low-power microscope paired with a digital camera. It sends a ring of white flash, LED, or laser light through the outer edge of your pupil while collecting the returning picture purely through the center, eliminating reflections off the front of your eye.
    \item \textbf{True-Color Tissue Contrast:} Visible white light interacts naturally with structures at the back of the eye: blood vessels absorb green light so they stand out crisply, central macular pigment absorbs blue light, and deeper pigment layers reflect warm orange-red light to create lifelike color photos.
    \item \textbf{Standard Versus Ultra-Widefield Views:} Based on the optics of monocular indirect ophthalmoscopy, regular fundus cameras deliver an upright, magnified view of the central retina---typically 2.5x magnification across 30--50 degrees, or up to 5x magnification across a zoomed 15-degree close-up. Ultra-widefield scanning laser cameras use curved mirrors to sweep gentle laser light across 140--200 degrees---often working well even without dilating drops.
    \item \textbf{Digital Capture, Stereo Pairs, and Special Filters:} A high-detail digital sensor records the reflected light as full-color photos, and two slightly offset shots can be paired as stereo images to give doctors a three-dimensional view of the optic nerve or macula. Special optical filters can also be switched on, such as green ``red-free'' light to highlight nerve fiber layers or autofluorescence mode to check the health of the retina's pigment layer.
\end{itemize}

\section{Imaging Mechanism}

\begin{enumerate}
    \item \textbf{Getting Lined Up:} You rest your chin and forehead comfortably against the camera frame and look steadily at a small glowing target light inside while the photographer lines up your pupil (often using dim invisible infrared light on non-dilating cameras) and focuses sharply on your retina.
    \item \textbf{Brief Flash or Gentle Laser Sweep:} When the photographer presses the button, a millisecond flash---or a quick gentle scanning laser sweep on ultra-widefield cameras---lights up the back of your eye without glare.
    \item \textbf{Capturing the Full-Color Photo:} Light reflecting off your retina travels back through the center of your pupil onto a sensitive digital camera sensor, converting the view into a sharp full-color digital image.
    \item \textbf{Balancing Color and Stitching Wide Views:} Built-in software balances color, trims edge glare, adds special filter views if requested, and can stitch multiple angles together into a wide panoramic map.
    \item \textbf{Instant Check and Permanent Record:} The crisp photo pops up right away so the photographer can confirm sharp focus, then saves permanently to your chart for side-by-side comparison at future visits, remote specialist review, or AI-assisted diabetic screening.
\end{enumerate}

\section{Key Technical Parameters}

\noindent\begin{tabularx}{\textwidth}{@{}l X@{}}
\toprule
\textbf{Parameter} & \textbf{Typical Value / Description} \\
\midrule
Physical Principle & Glare-free visible-light photography of the retina and optic nerve through the pupil \\
Energy / Signal Source & Visible white flash, LED, or multi-color scanning laser light \\
Typical Frequency / Energy & Ordinary visible light (400--700 nm); optional green red-free or autofluorescence filters \\
Spatial Resolution & Very fine 5--15 micrometer detail across standard 30--50 degree views (up to 140--200 degrees on ultra-widefield cameras) \\
Typical Scan Duration & 5--10 minutes for both eyes (plus 20--30 minutes beforehand if dilating eye drops are used) \\
Contrast Agent & None (completely noninvasive, no dye or injections) \\
Ionizing Radiation & None (uses ordinary visible light; zero X-rays) \\
\bottomrule
\end{tabularx}

\section{Common Symptoms That Prompt This Scan}
\begin{itemize}
    \item \textbf{Blurry Vision, Dark Spots, or Distortion:} Difficulty reading fine print, patchy vision, or new visual changes prompting a detailed color photograph of the retina and optic nerve.
    \item \textbf{Routine Diabetic or High Blood Pressure Eye Checks (No Symptoms):} Early diabetic eye disease and blood-pressure vessel damage rarely cause symptoms at first, making annual fundus photos vital.
    \item \textbf{Gradual Side-Vision Loss or High Eye Pressure:} Documenting the exact shape and color of the optic nerve over time to watch for glaucoma.
    \item \textbf{Tracking a Freckle or Spot Inside the Eye:} Baseline and follow-up photography of harmless retinal freckles (nevi) or macular changes to catch any subtle growth.
    \item \textbf{Persistent Headaches, Blood-Pressure Emergencies, or Medication Monitoring:} Checking urgently for swollen optic discs (papilledema) in patients with persistent headaches or severely elevated blood pressure, and screening the retina during long-term medications such as anti-malarial therapy.
\end{itemize}

\section{Clinical Usages}

\subsection{Primary Usages}

\begin{itemize}
    \item \textbf{Diabetic Eye Screening:} Quickly checks for tiny blood-vessel bulges (microaneurysms), small retinal bleeds, yellow fluid deposits, and abnormal new vessels---widely used in primary-care teleophthalmology and AI-assisted screening programs.
    \item \textbf{Tracking Glaucoma:} Creates a permanent visual record of the optic nerve's central cup, rim thickness, and tiny disc bleeds so doctors can spot subtle changes over years.
    \item \textbf{Monitoring Macular Degeneration (AMD):} Documents yellow aging deposits (drusen), pigment changes, worn-out thinning patches (geographic atrophy), and bleeding to track disease stage and treatment progress.
    \item \textbf{Checking Blood-Vessel and Blood-Pressure Effects:} Records blocked retinal veins or arteries, narrowed arteries from high blood pressure, fluffy white oxygen-starved spots (cotton-wool spots), and optic nerve swelling.
    \item \textbf{Watching Eye Freckles and Tumors:} Side-by-side serial photos track harmless choroidal freckles (nevi) or eye tumors closely to catch even tiny signs of growth or fluid change early.
\end{itemize}

\subsection{Secondary Usages}

\begin{itemize}
    \item \textbf{Children's Retinal Screening:} Specialized widefield digital cameras help screen premature infants for retinopathy of prematurity (ROP) and document retinal bleeding in pediatric eye exams.
    \item \textbf{Eye Infections and Inherited Conditions:} Records baseline and healing changes in eye inflammation (uveitis), viral or parasitic retinal infections, and inherited retinal conditions such as retinitis pigmentosa.
    \item \textbf{Retinal Tears and Neuro-Ophthalmic Evaluation:} Helps document retinal tears or detachments and supports neuro-ophthalmologists evaluating unexplained vision loss, poor optic-nerve blood flow (ischemic optic neuropathy), or optic disc swelling from raised brain pressure.
    \item \textbf{Special Filter Photography:} Red-free (green-filtered) photos sharpen nerve-fiber thinning and tiny vessel changes, while autofluorescence maps pigment wear and tear.
    \item \textbf{Research, Consultations, and Patient Education:} Provides standardized, easy-to-share images for clinical trials, second opinions, and showing patients their own eye health.
\end{itemize}

\section{Advantages}

\begin{itemize}
    \item \textbf{Fast, Painless, and Touch-Free:} Captures crisp retinal and optic nerve photos in minutes with zero radiation, no dye injections, and nothing touching your eye.
    \item \textbf{Permanent Side-by-Side Tracking:} Saves exact baseline pictures so your eye doctor can spot tiny changes in glaucoma, diabetes, macular degeneration, or eye freckles years later.
    \item \textbf{Ideal for Community Screening and AI:} Non-dilating and portable cameras make large-scale diabetic eye screening, remote specialist grading, and FDA-cleared AI screening easy and accessible.
    \item \textbf{Wide View of the Outer Retina:} Ultra-widefield cameras capture up to 140--200 degrees in a single click, spotting outer-retina problems that standard narrow views can miss.
    \item \textbf{Safe and Affordable:} Widely available, low-cost, and comfortable for adults and children alike.
\end{itemize}

\section{Disadvantages}

\begin{itemize}
    \item \textbf{Blurred by Eye Cloudiness:} Cloudy corneas, thick cataracts, floaters, or bleeding inside the eye scatter the camera light and can blur or block the retina.
    \item \textbf{Flat Two-Dimensional Picture:} Cannot measure cross-sectional thickness or tiny pockets of retinal swelling on its own, so doctors often pair it with an OCT scan.
    \item \textbf{Does Not Show Live Blood Flow:} Captures a still snapshot rather than watching active dye leakage or live circulation seen on fluorescein or ICG angiography.
    \item \textbf{Small Pupils and Eyelash Shadows:} Very small pupils can darken the edges of the photo, and eyelids or lashes can cast shadows on ultra-widefield shots, so dilating drops are sometimes needed for the clearest outer-retina view.
    \item \textbf{Mild Side Effects of Dilating Drops:} Photography itself has almost no risk beyond a brief bright flash; when dilating drops are used, blurry near vision and glare last several hours, and there is a small risk of triggering sudden high eye pressure in narrow-angle eyes.
\end{itemize}

\section{Contraindications}

\subsection{Absolute}

\begin{itemize}
    \item \textbf{None for Standard Non-Dilating Photos:} There are no strict medical reasons preventing non-mydriatic fundus photography, which is completely noninvasive and radiation-free.
\end{itemize}

\subsection{Relative}

\begin{itemize}
    \item \textbf{When Dilating Drops Should Be Avoided:} Very narrow drainage angles in the front of the eye or a history of untreated acute angle-closure glaucoma mean dilating drops are generally avoided in favor of non-dilating cameras.
    \item \textbf{Temporary Eye or Positioning Barriers:} Active contagious eye infections (like pink eye) call for waiting or strict isolation cleaning; extreme light sensitivity, involuntary eye movements (nystagmus) or eyelid squeezing, trouble sitting at the chin rest, or dense cataracts and eye bleeding can make clear photos difficult.
\end{itemize}

\section{Patient Safety Guidelines}

\begin{itemize}
    \item \textbf{Checking Eye Angles Before Dilation:} Before using dilating drops, your eye care team checks the front chamber of your eye and asks about narrow-angle glaucoma, choosing non-dilating photography whenever dilation is unsafe or unnecessary.
    \item \textbf{Clean Surfaces Between Every Patient:} Chin rests, forehead bars, and any contact surfaces are wiped with approved disinfectants between every patient, and non-urgent photos are postponed if you have an active contagious eye infection.
    \item \textbf{Safe Flash Brightness:} Cameras operate well within strict ANSI and ISO light-safety standards, keeping flash repeats minimal and pausing briefly for comfort, especially in light-sensitive patients or children.
    \item \textbf{Post-Dilation Safety Advice:} If dilating drops are used, wear sunglasses, avoid driving until blurry near vision and glare clear (usually 4--6 hours), and seek urgent eye care for severe eye pain, headache, rainbow halos around lights, or nausea.
\end{itemize}

\section{Prerequisite Tests}

\begin{itemize}
    \item \textbf{No Blood Tests or Prep Procedures:} Fundus photography is completely noninvasive and needs no blood work, IV lines, or medical clearance.
    \item \textbf{Quick Vision and Health Check:} A brief vision test and review of your eye and health history (such as diabetes, glaucoma, or past eye surgery) help guide which views to capture.
    \item \textbf{Front-of-Eye Check (If Dilating):} When dilating drops are planned, your doctor checks the front eye angle (and eye pressure when appropriate) to make sure dilation is safe.
    \item \textbf{Checking Eye Clarity:} A quick look ensures there are no dense cataracts blocking the lens or active contagious eye infections before photography.
\end{itemize}

\section{Preparing for the Scan}

\begin{itemize}
    \item \textbf{Contacts, Glasses, and Sunglasses:} Remove contact lenses if dilating eye drops are planned, and bring your regular glasses plus dark sunglasses to wear afterward for bright outdoor glare.
    \item \textbf{Ask About Dilation and Plan a Ride:} Check ahead whether your clinic uses a non-dilating camera or dilating drops. If dilation is planned, allow 20--30 minutes for drops to work and arrange a driver home, since close-up vision stays blurry for 4--6 hours.
    \item \textbf{Normal Meals and Eye Drops:} No fasting is needed---eat normally and take all your regular pills and prescribed glaucoma or eye drops on schedule unless your eye doctor says otherwise.
    \item \textbf{Simple Comfort Tips:} Skip heavy mascara or eye makeup that can flake or cast lash shadows on widefield photos, and be ready to rest your chin and forehead gently against the clean camera frame during brief bright flashes.
\end{itemize}

\section{During Scan}

\begin{itemize}
    \item \textbf{Sitting Comfortably and Looking Steady:} You will sit upright in front of the camera with your chin on the rest and forehead gently against the bar. Open your eyes wide when asked and focus steadily on the glowing target light inside.
    \item \textbf{Natural Blinking and Breathing:} Breathe normally and blink freely between shots to keep your eyes comfortable, pausing your blink for just a second when the photographer counts down to snap each photo.
    \item \textbf{What You Will Feel and See:} The camera never touches your eye and the test is completely painless. You will see a quick, bright flash for each photo (usually 1--3 shots per eye), leaving a harmless temporary colored afterimage that fades within a few minutes.
    \item \textbf{Speaking Up Anytime:} Tell the photographer if you need a quick pause to blink, adjust the chair, or rest your neck, or if dilating drops make your eyes sting or feel lightheaded.
    \item \textbf{How Long It Takes:} Taking the actual pictures takes just 5--10 minutes for both eyes, plus 20--30 minutes beforehand if dilating eye drops are used.
\end{itemize}

\section{Normal Appearance}

\begin{itemize}
    \item \textbf{Healthy Optic Nerve (Optic Disc):} A neat, pink-orange oval disc with sharp borders, a full healthy rim around the edge, and a small, symmetric central cup (typically 0.3--0.5 of the disc width).
    \item \textbf{Smooth Center of Vision (Macula and Fovea):} A smooth, slightly darker oval zone beside the optic nerve with a crisp tiny central shine (especially in younger eyes) and no yellow deposits, bleeding, or fluid spots.
    \item \textbf{Healthy Retinal Blood Vessels:} Smoothly branching arteries and veins emerging neatly from the optic nerve in balanced width, with no pinching, swelling, or abnormal lacy new vessels.
    \item \textbf{Even Retinal Background:} A smooth reddish-orange background (naturally lighter or darker depending on skin and eye pigment) free of bleeds, fluffy white spots, scars, or suspicious dark lumps.
\end{itemize}

\section{Observing for Abnormal Patterns}

\begin{itemize}
    \item \textbf{Diabetic and High Blood-Pressure Changes:} Tiny pinpoint red dots (microaneurysms), small blot or flame-shaped bleeds, sharp yellow fat deposits (hard exudates), fluffy white spots (cotton-wool spots), uneven veins, and delicate abnormal new vessels help grade diabetic and blood-pressure eye disease.
    \item \textbf{Glaucoma Optic Nerve Changes:} A deepened, vertically stretched optic cup, thin notched edges along the top or bottom rim, shifted vessels, tiny splinter bleeds on the disc edge, and dark wedge-shaped nerve-fiber shadows signal progressive glaucoma.
    \item \textbf{Macular Degeneration Signs:} Clusters of yellow spots under the macula (drusen), patchy pigment changes, sharply outlined pale thinning patches (geographic atrophy), or fresh bleeding and gray-green patches help tell dry from wet AMD.
    \item \textbf{Blocked Vessels and Nerve Swelling:} Widespread flame-shaped bleeds and swollen twisted veins across the retina point to a blocked central vein, while a pale retina with a bright red center spot (``cherry-red spot'') signals a blocked central artery. Swollen, blurry optic nerve borders in both eyes suggest raised pressure around the brain (papilledema).
    \item \textbf{Harmless Look-Alikes:} Small pupils can leave dark crescents around the edges, cataracts create fuzzy glare halos, and eyelids or lashes can cast dark shadows on widefield photos. Harmless feathery white nerve coatings or tiny sparkling vitreous specks (asteroid hyalosis) can also resemble cotton-wool spots or exudates at first glance.
\end{itemize}

\section{Dealing with Uncertainties}

\begin{itemize}
    \item \textbf{Cloudy Views and Shadows:} Thick cataracts, corneal haze, eye bleeding, tiny undilated pupils, or eyelash shadows can hide retinal spots or look like dark marks. Using dilating drops, gently holding the eyelid, and retaking aligned views often clears the picture.
    \item \textbf{No Cross-Section Depth:} Because color photos are flat, subtle swelling in the macula, shallow fluid under the retina, early wet AMD, or mild optic nerve swelling cannot be measured or ruled out from photos alone.
    \item \textbf{Harmless Spots That Mimic Disease:} Benign eye freckles can look similar to small eye melanomas, harmless buried calcium spots (optic disc drusen) can look like nerve swelling, and benign white nerve fibers can mimic cotton-wool spots.
    \item \textbf{Adding Other Eye Scans:} Doctors clear up uncertain photos by pairing them with cross-sectional OCT scans, autofluorescence photos, gentle eye ultrasound (for dark spots), dye angiography (for suspected leakage), or comparing against your past baseline photos.
\end{itemize}

\section{Augmenting Methods}

\begin{itemize}
    \item \textbf{OCT and Dye-Free OCT Angiography (OCTA):} Cross-sectional OCT measures exact retinal swelling and glaucoma nerve-fiber thickness, while OCTA maps tiny blood vessels without any dye injection.
    \item \textbf{Autofluorescence and Special Color Filters:} Autofluorescence highlights pigment wear and tear and dry AMD borders, while green red-free and infrared filters sharpen nerve-fiber thinning and deep spots.
    \item \textbf{Fluorescein and ICG Dye Angiography:} Confirm live blood-vessel leakage, outer-retina oxygen starvation, and deep abnormal vessels when still color photos are not enough.
    \item \textbf{Eye Ultrasound, Visual Field Tests, and AI Screening:} Ultrasound measures the thickness of dark spots inside the eye, visual field tests measure side-vision loss in glaucoma, and FDA-cleared AI tools help screen large numbers of diabetic eye photos accurately.
\end{itemize}

\section{Comparison with Alternative Modalities}

\begin{itemize}
    \item \textbf{Compared with OCT:} OCT measures microscopic cross-sectional thickness to track macular swelling and glaucoma nerve layers, while fundus photography gives true-color surface views ideal for diabetic screening and watching pigmented spots.
    \item \textbf{Compared with Dye Angiography (FFA/ICGA):} Dye angiography reveals active fluid leakage and live circulation, while fundus photography is faster, dye-free, completely noninvasive, and ideal for routine screening and baseline records.
    \item \textbf{Compared with a Standard Dilated Eye Exam:} Looking through a slit lamp gives your eye doctor a live 3D view, while fundus photography creates a permanent, objective picture for future comparison, remote care, and AI screening.
\end{itemize}

\section{Follow-up Steps}

\begin{itemize}
    \item \textbf{After-Exam Care:} If no dilating drops were used, you can return to everything immediately. If your eyes were dilated, wear sunglasses and avoid driving until glare and near-blur fade (4--6 hours), and seek urgent care for severe eye pain, halos around lights, or nausea.
    \item \textbf{Prompt Results and Referrals:} Your baseline photos are saved for future comparison. If screening spots sight-threatening diabetic eye disease, optic nerve swelling, a blocked vessel, or a suspicious dark spot, you will be referred promptly to an eye specialist.
    \item \textbf{Extra Eye Tests if Needed:} Abnormal or unclear photos are often followed by an OCT scan for retinal swelling or glaucoma, visual field testing, eye ultrasound for dark spots, or dye angiography for suspected leakage.
    \item \textbf{Routine Checkup Schedule:} Repeat photo intervals depend on your eye health---for example, once a year for mild or no diabetic eye disease, every 3--6 months for moderate diabetic changes or closely watched eye freckles, and periodic optic nerve photos for glaucoma tracking.
\end{itemize}

\section{Key References}

\begin{itemize}
    \item American Academy of Ophthalmology (AAO). \textit{Diabetic Retinopathy Preferred Practice Pattern}. \textit{Ophthalmology}, 2019 (updated 2024).
    \item American Academy of Ophthalmology (AAO). \textit{Primary Open-Angle Glaucoma Preferred Practice Pattern}. \textit{Ophthalmology}, 2021.
    \item Abramoff MD, Lavin PT, Birch M, Shah N, Folk JC. Pivotal Trial of an Autonomous AI-Based Diagnostic System for Detection of Diabetic Retinopathy in Primary Care Offices. \textit{npj Digital Medicine}, 2018;1:39.
    \item Panwar N, Huang P, Lee J, et al. Fundus Photography in the 21st Century---A Review of Recent Technological Advances and Their Implications for Worldwide Healthcare. \textit{Telemedicine and e-Health}, 2016;22(3):198--208.
    \item Mishra C, Tripathy K. \textit{Fundus Camera}. StatPearls, National Library of Medicine (NIH). Available at: \url{https://www.ncbi.nlm.nih.gov/books/NBK585111/}.
    \item Penn Medicine. \textit{Fundus Photography}. University of Pennsylvania Health System. Available at: \url{https://www.pennmedicine.org/treatments/fundus-photography}.
\end{itemize}

\section{Further Reading}
\begin{itemize}
  \item Wikipedia contributors. \textit{Fundus photography}. Wikipedia, The Free Encyclopedia. Available at: \url{https://en.wikipedia.org/wiki/Fundus_photography}.
\end{itemize}
